
Benchmark contamination---the presence of test data in training corpora---is documented in large language model evaluation, yet no prior work quantifies how contamination translates to score inflation. This paper proposes a checkpoint-gradient methodology for measuring contamination-inflation correlation: by comparing model checkpoints at different training stages with varying contamination exposure, correlation analysis becomes feasible without contamination-free baselines. Capability detrending via out-of-distribution perplexity isolates contamination effects from legitimate capability gains. In proof-of-concept validation using simulated data across 80 Pythia checkpoint-benchmark pairs (5 model sizes $\times$ 4 benchmarks $\times$ 4 checkpoints), the methodology yields a positive correlation (Spearman r = 0.326, p = 0.003) between contamination proxy and inflation residual. Per-benchmark analysis shows MMLU (r = 0.55, p = 0.01) and ARC-Challenge (r = 0.53, p = 0.02) with statistically significant correlations, while HellaSwag (r = 0.30, p = 0.20) and WinoGrande (r = 0.25, p = 0.29) do not reach significance. Separate n-gram overlap analysis on a 50,000-document Pile subset confirms detection mechanism validity, identifying individual MMLU items with up to 17.4\% 13-gram overlap. These preliminary findings suggest contamination-inflation correlation may be measurable pending full-scale validation with real Pythia checkpoint evaluations, which requires approximately 144 GPU-hours.

